\documentclass[a4paper,11pt]{article}
\usepackage[utf8]{inputenc}
\usepackage[T1]{fontenc}
\usepackage[french]{babel}
\usepackage[top=2cm, bottom=2cm, left=2.3cm, right=2.3cm]{geometry}
\usepackage{xcolor}
\usepackage{booktabs}
\usepackage{tabularx}
\usepackage{enumitem}
\usepackage{listings}
\usepackage{float}
\usepackage{amssymb}
\usepackage{tcolorbox}
\usepackage{hyperref}
\hypersetup{colorlinks=true, linkcolor=primary, urlcolor=primary, citecolor=primary}

\definecolor{primary}{HTML}{1A237E}
\definecolor{success}{HTML}{2E7D32}
\definecolor{warning}{HTML}{E65100}
\definecolor{lightblue}{HTML}{E3F2FD}
\definecolor{lightorange}{HTML}{FFF3E0}
\definecolor{codebg}{HTML}{F5F5F5}
\definecolor{codeblue}{rgb}{0.13,0.29,0.53}
\definecolor{codegreen}{rgb}{0.0,0.50,0.0}

\lstdefinestyle{sh}{backgroundcolor=\color{codebg}, basicstyle=\ttfamily\footnotesize,
  keywordstyle=\color{codeblue}\bfseries, commentstyle=\color{codegreen}\itshape,
  breaklines=true, frame=single, showstringspaces=false, tabsize=2, columns=fullflexible}
\lstset{style=sh}

\title{\textbf{Perception LiDAR, navigation autonome (SLAM / Nav2)\\ et communication V2V}\\[2mm]
\large Rapport technique hebdomadaire --- ProtoVA2 (Jetson Orin Nano, ROS\,2 Humble)}
\author{Douae Sebti --- ALTEN Labs}
\date{Semaine du 23 au 26 juin 2026}

\begin{document}
\maketitle

\begin{abstract}
\noindent Ce rapport présente les travaux que j'ai menés cette semaine sur le
prototype ProtoVA2, sur trois axes : (i) la mise en service du \textbf{LiDAR} et la
chaîne de perception, (ii) la \textbf{navigation autonome} --- d'abord une approche
réactive, puis une architecture \emph{cartographiée} (SLAM + \texttt{Nav2}) ---, et
(iii) la \textbf{communication V2V} (\emph{vehicle-to-vehicle}) au standard
européen ETSI, via la pile \texttt{Vanetza}. Pour chaque axe, je détaille les choix
techniques, les résultats obtenus, les limites identifiées et les perspectives pour
la semaine suivante.
\end{abstract}

\tableofcontents
\bigskip\hrule

% =====================================================================
\section{Contexte et objectifs}
ProtoVA2 est un véhicule à l'échelle (direction \textbf{Ackermann}, propulsion par
moteur brushless/ESC, calculateur Jetson Orin Nano sous ROS\,2 Humble). Cette
semaine, je me suis fixé deux objectifs : doter le véhicule d'une capacité de
\textbf{navigation autonome} fondée sur la perception LiDAR, et mettre en place une
\textbf{communication V2V} entre véhicules. J'ai mené ces deux volets en parallèle.

% =====================================================================
\section{Perception LiDAR}

\subsection{Matériel et intégration}
Le capteur est un \textbf{RPLIDAR A1} (firmware 1.29), relié par un pont USB-série
Silicon Labs \textbf{CP2102} (\texttt{10c4:ea60}) exposé sur \texttt{/dev/ttyUSB0},
à \textbf{115200} bauds. Le pilote ROS\,2 \texttt{rplidar\_ros} publie le flux
\texttt{/scan} (\texttt{sensor\_msgs/LaserScan}).

\paragraph{Mise en service --- diagnostic.} Le capteur ne renvoyait initialement
aucune donnée (\texttt{SL\_RESULT\_OPERATION\_TIMEOUT}) bien que le moteur tourne.
J'ai mené un diagnostic par élimination (débit série, droits d'accès, interférence
ModemManager, sonde série bas niveau, lignes DTR/RTS) qui m'a permis d'écarter
toute cause logicielle et d'identifier un \textbf{défaut de connecteur de données} ;
le ré-emboîtement a rétabli la communication.

\subsection{Caractéristiques du flux \texttt{/scan}}
\begin{center}
\begin{tabular}{@{}ll@{}}
\toprule
Paramètre & Valeur \\
\midrule
Fréquence de scan & $\approx$ 7,4 Hz \\
Points par tour & $\approx$ 1080 \\
Portée & 0,15 -- 12 m \\
Repère (\texttt{frame\_id}) & \texttt{laser} \\
Mode & \emph{Sensitivity}, 8 kHz \\
\bottomrule
\end{tabular}
\end{center}

\paragraph{Montage et repère.} Le LiDAR est monté « moteur vers l'arrière » :
l'angle nul du scan pointe vers l'arrière du véhicule. Ce décalage de 180\textdegree{}
est encodé directement dans la description URDF (articulation
\texttt{laser\_joint}, \texttt{rpy = 0\ 0\ $\pi$}), de sorte que la chaîne de
transformations (TF) restitue automatiquement l'orientation correcte.

\subsection{Visualisation (RViz)}
J'ai visualisé le flux sur le PC via WiFi (DDS natif). En fixant le \emph{Fixed
Frame} sur \texttt{laser}, l'affichage du nuage de points ne requiert pas d'arbre
TF complet, ce qui m'a permis de valider rapidement la qualité du scan.

% =====================================================================
\section{Première approche : navigation réactive}
J'ai d'abord implémenté une navigation \textbf{réactive} de type
\emph{follow-the-gap} (paquet ROS\,2 \texttt{navigation}, n{\oe}ud
\texttt{follow\_the\_gap}).

\paragraph{Principe.} À chaque scan : je sélectionne le secteur avant
($\pm$90\textdegree), j'étends les obstacles à la demi-largeur du véhicule
(\emph{disparity extender}), je place une bulle de sécurité autour du point le plus
proche, je cherche le plus grand espace libre, puis je vise le centre de sa zone la
plus dégagée. La sortie est une consigne \texttt{geometry\_msgs/Twist}
(\texttt{angular.z} = angle de braquage servo $\in[28;138]$, centre 83 ;
\texttt{linear.x} = vitesse en m/s), publiée sur \texttt{/cmd\_vel\_auto} et
arbitrée par \texttt{twist\_mux} en priorité basse (sous l'AEB et la reprise pilote).

\paragraph{Résultats.} J'ai validé le n{\oe}ud sur scan synthétique (champ libre,
mur frontal, obstacle latéral) puis sur données réelles ; j'ai vérifié la chaîne
complète capteur $\rightarrow$ décision $\rightarrow$ actionneur (Pico).

\begin{tcolorbox}[colback=lightorange, colframe=warning, title=\textbf{Limites de l'approche réactive}]
La navigation réactive s'est révélée \textbf{insuffisante} : oscillations de la
consigne de braquage (la « plus grande ouverture » alterne entre côtés) et
évitement d'obstacles peu robuste car \emph{sans mémoire} de l'environnement (le
véhicule ne sait pas où il est, il ne fait que réagir à l'instant). Ces limites
m'ont conduite à passer à une \textbf{navigation cartographiée}.
\end{tcolorbox}

% =====================================================================
\section{Navigation cartographiée : SLAM et Nav2}
J'ai retenu la navigation fondée sur une carte (standard de l'état de l'art
robotique), qui apporte planification globale et évitement d'obstacles par
\emph{costmap}.

\subsection{Chaîne de transformations (TF) et odométrie}
La cohérence spatiale repose sur l'arbre TF :
\[
\texttt{map} \rightarrow \texttt{odom} \rightarrow \texttt{base\_link}
\rightarrow \texttt{laser}
\]
\begin{itemize}[leftmargin=1.4em, nosep]
  \item \texttt{base\_link} $\rightarrow$ \texttt{laser} : statique, issu de l'URDF
        (\texttt{car\_description} + \texttt{robot\_state\_publisher}).
  \item \texttt{odom} $\rightarrow$ \texttt{base\_link} : \textbf{odométrie}. Le
        paquet \texttt{localisation} fournit une odométrie laser (\texttt{rf2o}),
        une odométrie Ackermann (encodeur + angle de braquage) et une fusion
        \textbf{EKF} (\texttt{robot\_localization}). Pour la cartographie, j'ai
        utilisé \texttt{rf2o} seule afin de limiter les dépendances.
  \item \texttt{map} $\rightarrow$ \texttt{odom} : fourni par le SLAM.
\end{itemize}

\subsection{SLAM (cartographie)}
La cartographie s'appuie sur \texttt{slam\_toolbox} (mode asynchrone en ligne).
J'ai rendu opérationnelle la chaîne LiDAR + TF + odométrie laser +
\texttt{slam\_toolbox} : l'arbre TF est complet et la grille d'occupation
(\texttt{/map}) se construit. Je sauvegarde la carte au format ROS standard
(\texttt{.pgm} + \texttt{.yaml}) via \texttt{nav2\_map\_server}.

\subsection{Pile de navigation Nav2}
J'ai installé la pile \texttt{Nav2} (\texttt{navigation2}, \texttt{nav2\_bringup}).
Composants clés retenus :
\begin{itemize}[leftmargin=1.4em, nosep]
  \item \textbf{Contrôleur} : \texttt{RegulatedPurePursuitController}, adapté à une
        cinématique \textbf{Ackermann} (suivi de trajectoire par point de visée).
  \item \textbf{Costmaps} en couches : couche \emph{statique} (la carte = murs
        permanents), couche \emph{obstacles} (LiDAR temps réel), couche
        \emph{inflation} (marge de sécurité).
  \item \textbf{Localisation} : \texttt{AMCL} sur la carte enregistrée.
\end{itemize}

\begin{tcolorbox}[colback=lightblue, colframe=primary, title=\textbf{Gestion des obstacles dynamiques (mobilier)}]
La problématique des obstacles mobiles (tables, chaises déplacées) est traitée
nativement par l'architecture en couches de Nav2 : la \textbf{carte statique} ne
contient que la structure permanente (murs), tandis que la \textbf{couche obstacles}
intègre en temps réel ce que voit le LiDAR et \emph{efface} l'espace redevenu libre.
Le planificateur contourne donc les obstacles \emph{présents}, indépendamment de
l'état de la carte au moment de sa création.
\end{tcolorbox}

% =====================================================================
\section{Communication V2V (Vanetza, ETSI CAM)}
En parallèle de la navigation, j'ai mis en place une communication
\textbf{véhicule-à-véhicule} au standard européen.

\subsection{Standard et pile logicielle}
J'ai retenu le standard \textbf{ETSI C-ITS} (messages \textbf{CAM}, \emph{Cooperative
Awareness Message}, et à terme \textbf{DENM} pour les événements) --- et non le BSM
américain, puisque nous sommes en Europe. J'utilise la pile open-source
\textbf{Vanetza}, qui implémente la pile ITS-G5 (GeoNetworking, BTP, CAM/DENM en
ASN.1). Ne disposant pas de radio 802.11p, j'exploite l'application
\texttt{socktap} de Vanetza sur le \textbf{WiFi} (lien UDP multicast) : les messages
échangés sont de vrais CAM ETSI, seule la couche radio est remplacée.

\subsection{Échange de CAM}
J'ai compilé Vanetza (option \texttt{socktap}, backend OpenSSL) sur le PC et la
Jetson, puis lancé \texttt{socktap} de part et d'autre. \textbf{Résultat :} échange
de CAM \textbf{bidirectionnel} confirmé sur WiFi (chaque véhicule reçoit le
\texttt{Station ID} de l'autre), sans matériel dédié ni privilèges \texttt{root}.
Chaque CAM ($\approx$99 octets, $\approx$1 Hz) contient l'en-tête ITS (version,
type, identifiant station) et le conteneur \emph{CoopAwareness} (position GPS,
cap, vitesse, type et dimensions du véhicule).

\subsection{Intégration ROS2 et alerte de proximité}
J'ai développé un n{\oe}ud ROS\,2 (\texttt{v2v\_node.py}) qui exploite les CAM
reçus : il en extrait la position des autres véhicules, calcule la distance et
publie \texttt{/v2v/remote\_vehicles} ainsi qu'une \texttt{/v2v/alert} en cas de
proximité. \textbf{Résultat :} véhicule distant à 5\,m $\rightarrow$ alerte ; à
1,1\,km $\rightarrow$ pas d'alerte (la logique discrimine correctement).

\subsection{Position dynamique issue de l'odométrie}
Pour que le CAM porte la position \emph{réelle} du véhicule, j'ai ajouté à
\texttt{socktap} un fournisseur de position alimenté par UDP, puis un n{\oe}ud
(\texttt{odom\_to\_gps}) qui convertit l'odométrie locale (mètres) en coordonnées
GPS. \textbf{Résultat :} en poussant le véhicule (aller-retour d'$\approx$1,7\,m),
la distance mesurée par l'autre véhicule a suivi en direct
($3{,}3 \rightarrow 5{,}3 \rightarrow 3{,}4$ m), l'alerte disparaissant au-delà de
5\,m puis revenant --- le mouvement physique réel voyage bien dans les CAM.

\subsection{Vérification de la communication}
J'ai prévu trois niveaux de vérification : l'affichage applicatif de
\texttt{socktap} (\texttt{--print-rx-cam}), l'observation des trames sur le réseau
avec \texttt{tcpdump} (multicast \texttt{239.118.122.97}, port \texttt{8947}), et
le décodage complet des champs CAM sous \textbf{Wireshark}.

% =====================================================================
\section{Limites actuelles}
\begin{itemize}[leftmargin=1.4em]
  \item \textbf{Carte du lieu (navigation).} La carte existante fournie est une
        image (PNG) sans métadonnées (\texttt{.yaml} : résolution, origine) ; elle
        n'est pas directement exploitable par Nav2. Une cartographie propre par SLAM
        est nécessaire.
  \item \textbf{Interface de commande Ackermann.} Nav2 produit une consigne
        \texttt{Twist} (vitesse linéaire, vitesse de rotation), alors que le
        véhicule attend un \textbf{angle de braquage servo}. Un n{\oe}ud de
        conversion \emph{twist $\rightarrow$ angle de braquage}
        ($\delta = \arctan(L\,\omega / v)$, $L$ empattement) reste à écrire.
  \item \textbf{Transport des grosses données sur WiFi.} Le topic \texttt{/map}
        (et les costmaps) ne traverse pas le WiFi avec la configuration FastDDS par
        défaut ; un profil adapté (ou RViz exécuté localement) sera requis.
  \item \textbf{V2V.} La communication se limite pour l'instant aux \textbf{CAM} et
        à une alerte de proximité ; les \textbf{DENM} (événements) et la
        \textbf{sécurité} (messages signés par certificats) restent à ajouter.
  \item \textbf{Fiabilité mécanique du LiDAR.} Le connecteur de données se
        désolidarise lors des manipulations du véhicule ; une fixation pérenne est
        à prévoir.
\end{itemize}

% =====================================================================
\section{Perspectives --- semaine du 30 juin 2026}
\begin{enumerate}[leftmargin=1.6em]
  \item \textbf{Cartographie propre} du site par SLAM, avec sauvegarde \texttt{.pgm}
        + \texttt{.yaml} (résolution et origine correctes).
  \item \textbf{N{\oe}ud de conversion Ackermann} (twist Nav2 $\rightarrow$ angle de
        braquage servo $\rightarrow$ trame Pico) et intégration via \texttt{twist\_mux}.
  \item \textbf{Localisation AMCL} sur la carte, réglage des \textbf{costmaps} et du
        contrôleur Pure Pursuit, puis \textbf{navigation autonome de bout en bout}
        (objectif envoyé depuis RViz, évitement d'obstacles réels).
  \item \textbf{V2V} : ajout des \textbf{DENM} (alerte de freinage / obstacle),
        puis sécurité par certificats ; à terme, coupler l'alerte V2V au freinage
        d'urgence (AEB).
  \item \textbf{Transport WiFi} : profil FastDDS pour les topics volumineux.
\end{enumerate}

\paragraph{Synthèse de l'état d'avancement.}
\begin{center}
\begin{tabularx}{\textwidth}{@{}X l@{}}
\toprule
\textbf{Élément} & \textbf{État} \\
\midrule
Perception LiDAR (\texttt{/scan}, diagnostic, RViz) & \textcolor{success}{Opérationnel} \\
Navigation réactive (\texttt{follow\_the\_gap}) & \textcolor{success}{Réalisée} (abandonnée : limites) \\
SLAM (\texttt{slam\_toolbox}, TF, odométrie) & \textcolor{success}{Opérationnel} \\
Installation Nav2 & \textcolor{success}{Terminée} \\
V2V : échange CAM (Vanetza) + ROS2 + alerte & \textcolor{success}{Opérationnel} \\
V2V : position dynamique (odométrie réelle) & \textcolor{success}{Opérationnel} \\
Carte géoréférencée du site & \textcolor{warning}{À produire} \\
Conversion Ackermann (twist $\rightarrow$ servo) & \textcolor{warning}{À développer} \\
Navigation Nav2 de bout en bout & \textcolor{warning}{À réaliser} \\
V2V : DENM + sécurité & \textcolor{warning}{À développer} \\
\bottomrule
\end{tabularx}
\end{center}

\end{document}
